在处理具备高度复杂边界约束的多体非平滑碰撞验证中，如何基于低开销重构建稳健的时空融合接触探测是一个极具挑战的核心命题。
本文跳出了增加局部晶格分辨率的工程算力内耗盲区，深挖隐式体素的二阶泰勒展开推导空间。提出了一种无需拆解修改核心互补 LCP 底层闭源逻辑的高阶补偿注入方法：借由提取纳米等级网格体系内的宏距 Hessian 张量（即 Weingarten 映射度量），以补偿预测的形式抵御了切向滑动外流和尖锐角过度估计产生的巨大数值刺激反馈。

通过凸轮-从动件实验数据，我们充分证实了该方法对冲量畸变拥有无损削去的杰出能力。二阶算法在全局平滑流形上表现出极高的收敛性，并能够突破传统多体系统临界积分域的两至三倍大步长（$\Delta t = 0.01\text{s}$），完全消除了传统一阶 SDF 带来的穿透与数值崩溃问题。此时算法展现出了直接搭载嵌入工程级主干网的极具吸引力的前景。

然而，本文同样通过直齿轮复杂啮合的基准测试，严格标定了该高阶方法的数学局限性。研究表明，在遭遇几何刚硬边缘与不可导的拓扑奇点（$C^0$ 奇异特征）时，依赖局部导数的空间泰勒级数展开必然遭遇 Hessian 矩阵特征值爆炸，导致局部补偿退化为高频噪声注入。针对这一极具价值的发现，我们论证了任何依赖离散点切线的碰撞测度算法在极限界面的天然缺陷，并由此确立了接触力学算法未来的技术路线演进方向。

在未来的推进工作中，针对带有锐利棱角及极小间隙的复杂装配体，研究的范式应当从“点导数预测”（Point-based Differential Approximation）跨越至“体积积分 SDF”（Volumetric Integral SDF）框架。通过将穿透程度的度量从逐点梯度距离转化为域内的体能量积分，可利用连续介质力学的几何体积作为天然的低通滤波器（Low-pass Filter），从而彻底隔离 $C^0$ 奇异点带来的局部高频干扰。这一由导数范式走向积分范式的转换，将是全面解决非连续几何体隐式接触碰撞的终局拼图。